\documentclass[11pt]{article}
\usepackage[T1]{fontenc}
\usepackage{lmodern}
\usepackage{amsmath,amssymb,amsthm,mathtools}
\usepackage{enumitem}
\usepackage{graphicx}
\usepackage[margin=1in]{geometry}
\usepackage[colorlinks=true,linkcolor=blue,citecolor=blue,urlcolor=blue]{hyperref}
\setlength{\emergencystretch}{2em}

\newtheorem{theorem}{Theorem}
\newtheorem{lemma}{Lemma}
\newtheorem{corollary}{Corollary}
\newtheorem{remark}{Remark}
\newcommand{\R}{\mathbb{R}}
\newcommand{\B}{\mathcal{B}}
\newcommand{\A}{\mathcal{A}}
\newcommand{\N}{\mathcal{N}}
\newcommand{\Hh}{\mathcal{H}}
\newcommand{\Pcal}{\mathcal{P}}

\title{Question 6 Below the Critical Exponent:\\
Cantor-Product Embeddings Force Undecidable Semilocalizability}
\author{Run \texttt{fa\_banach\_001}, agent \texttt{agent\_lane\_09}}
\date{9 August 2026}

\begin{document}
\maketitle

\section*{Status and scope}

Candidate strong partial resolution, likely valid, pending specialist review.
Question~6 of De Pauw asks, for $0<d<1$ and $d\ne\frac12$, whether
\[
  ([0,1],\A_{\Hh^d},\Hh^d)
\]
is consistently not semilocalizable~\cite{DePauw}.  This packet answers the
question affirmatively for every $0<d<\frac12$.  Combined with the source
paper's theorem at $d=\frac12$, this settles the whole lower range
$0<d\leq\frac12$.  No claim is made for $\frac12<d<1$.

The new step is to place a bi-Lipschitz copy of the product $C_d\times C_d$
inside $[0,1]$ when $d<\frac12$.  De Pauw's vertical--horizontal theorem
then gives consistent nonlocalizability on this compact copy.  A trace and
principal-ideal argument transfers the obstruction to the whole interval.

\section*{Source question}

The question occurs in Section~11.6, printed and PDF page~30 of
arXiv:1909.10190v2.  The crop includes the explanation of the product
structure and the complete statement of Q6.

\begin{center}
\includegraphics[width=\linewidth]{figures/open_problem_crop.png}
\end{center}

Transcription: for $0<d<1$, $d\ne\frac12$, is the measure space
$([0,1],\A_{\Hh^d},\Hh^d)$ consistently not semilocalizable?

\section*{Definitions and source inputs}

For an outer measure $\mu$ on $X$, let $\A_\mu$ be its Caratheodory
$\sigma$-algebra and put
\[
 \A_\mu^f=\{F\in\A_\mu:\mu(F)<\infty\},\qquad
 \N_\mu^{\mathrm{loc}}
 =\{N\in\A_\mu:\mu(N\cap F)=0\text{ for every }F\in\A_\mu^f\}.
\]
A measurable space with negligibles $(X,\A,\N)$ is \emph{localizable} when
every family in $\A$ has an essential supremum modulo $\N$; equivalently,
the quotient Boolean algebra $\A/\N$ is order complete.  A measure space
$(X,\A,\mu)$ is \emph{semilocalizable} when
$(X,\A,\N_\mu^{\mathrm{loc}})$ is localizable~\cite{DePauw}.

We use two results proved in the source paper.

\begin{enumerate}[label=\textup{(D\arabic*)}]
\item\label{input:vh}
Under the relatively consistent inequality
\[
 \operatorname{non}(\N_{\mathcal L^1})
 <\operatorname{cov}(\N_{\mathcal L^1}),
\]
De Pauw's Theorem~8.3 says that $(X,\A,\N)$ is not localizable if it
admits two Polish diffuse probability parameter spaces $S,T$ and measurable
sets $V_s,H_t$ such that:
\begin{enumerate}[label=\textup{(\alph*)}]
\item $\varnothing\ne V_s\cap H_t\in\N$ for every $s,t$;
\item if $V_s\cap Z\in\N$, then the outer probability of
$\{t:H_t\cap V_s\cap Z\ne\varnothing\}$ is zero;
\item the symmetric assertion holds after exchanging $V,S$ with $H,T$.
\end{enumerate}
\item\label{input:ch}
Under the Continuum Hypothesis, every Borel-regular outer measure on a
Polish space has the almost-decomposition property used in the paper; this
implies semilocalizability (Propositions~5.3--5.4 of the source).
\end{enumerate}

\section*{Proof intuition}

Let $C_d$ be a two-branch Cantor set whose contraction ratio is
$r=2^{-1/d}$.  The identity $2r^d=1$ gives its Hausdorff dimension $d$.
When $d<\frac12$, one has $4r<1$, so four separated copies of ratio $r$
fit inside an interval.  Pair the two binary digits of a point of
$C_d\times C_d$ into one four-valued digit.  This produces a bi-Lipschitz
coding of the product onto a compact four-branch Cantor set $K_d\subset[0,1]$.

The images of the vertical and horizontal leaves are $d$-dimensional sets of
finite positive measure and intersect in exactly one point.  If a measurable
set is locally null on one leaf, it is genuinely $\Hh^d$-null there; the
opposite parameter incidence set is a Lipschitz image and is therefore null.
This is precisely the input required by \ref{input:vh}.

Finally, the measurable quotient for $K_d$ is the principal ideal below the
class of $K_d$ in the interval quotient.  Principal ideals of complete
Boolean algebras are complete, so localizability of the interval would force
localizability of $K_d$, a contradiction.

\section*{The Cantor-product embedding}

Fix $0<d<\frac12$ and set
\[
 r=2^{-1/d},\qquad g_C=1-2r,\qquad g_K=\frac{1-4r}{3}.
\]
Thus $0<r<\frac14$ and $g_C,g_K>0$.  Let $C_d$ be the attractor of
\[
 S_0(x)=rx,\qquad S_1(x)=1-r+rx,
\]
and let $K_d$ be the attractor of the four similarities
\[
 P_a(x)=rx+\frac{a(1-r)}3,\qquad a=0,1,2,3.
\]
The four first-generation intervals are disjoint, with adjacent gap $g_K$.

\begin{lemma}[Digit pairing is bi-Lipschitz]\label{lem:bilip}
There is a bi-Lipschitz bijection
\[
 \Phi:C_d\times C_d\longrightarrow K_d.
\]
With the Euclidean product metric it satisfies
\[
 \frac{g_K}{\sqrt2}\,\|(x,y)-(x',y')\|_2
 \leq |\Phi(x,y)-\Phi(x',y')|
 \leq \frac1{g_C}\,\|(x,y)-(x',y')\|_2.
\]
\end{lemma}

\begin{proof}
Strong separation gives every point of $C_d$ a unique binary address.  If
$x$ and $y$ have addresses $(i_n)$ and $(j_n)$, define $\Phi(x,y)$ to be the
point of $K_d$ with address $(2i_n+j_n)$.  This is a bijection because digit
pairing is a bijection from $\{0,1\}^2$ to $\{0,1,2,3\}$ at every level.

Suppose two paired addresses first differ at level $n$.  On $C_d$, two
binary addresses first differing at level $n$ have distance between
$g_Cr^{n-1}$ and $r^{n-1}$.  Consequently, on $C_d\times C_d$,
\[
 g_Cr^{n-1}\leq \|(x,y)-(x',y')\|_2
 \leq \sqrt2\,r^{n-1}.
\]
On $K_d$, two four-valued addresses first differing at level $n$ have
distance between $g_Kr^{n-1}$ and $r^{n-1}$.  Combining these two pairs of
bounds gives the displayed inequalities.
\end{proof}

The open set condition gives
\begin{equation}\label{eq:cmeasure}
 0<\Hh^d(C_d)<\infty,
\end{equation}
since $2r^d=1$.  Notice also that $\dim_H K_d=2d$; the point of the
construction is not that $K_d$ itself has finite $\Hh^d$ measure, but that
its two transverse leaf families do.

\section*{The vertical--horizontal obstruction on $K_d$}

For $s,t\in C_d$, define compact subsets of $K_d$ by
\[
 V_s=\Phi(\{s\}\times C_d),\qquad
 H_t=\Phi(C_d\times\{t\}).
\]

\begin{lemma}[The source theorem applies on $K_d$]\label{lem:vh}
Let $\phi=\Hh^d|_{K_d}$, $\A=\A_\phi$, and
$\N=\N_\phi^{\mathrm{loc}}$.  Under
$\operatorname{non}(\N_{\mathcal L^1})<
\operatorname{cov}(\N_{\mathcal L^1})$, the measurable space with
negligibles $(K_d,\A,\N)$ is not localizable.
\end{lemma}

\begin{proof}
By Lemma~\ref{lem:bilip} and \eqref{eq:cmeasure}, every $V_s$ and $H_t$ is
compact and has finite positive $\Hh^d$ measure.  Moreover
\[
 V_s\cap H_t=\{\Phi(s,t)\},
\]
which is nonempty and $\Hh^d$-null because $d>0$.

Give $C_d$ the normalized diffuse Borel probability
$\Hh^d|_{C_d}/\Hh^d(C_d)$.  Fix $Z\in\A$.  If $V_s\cap Z\in\N$, then the
finite-measure test set $V_s$ in the definition of local nullity yields
\[
 \Hh^d(V_s\cap Z)=0.
\]
The map $V_s\to C_d$, $\Phi(s,t)\mapsto t$, is Lipschitz: it is a coordinate
projection after the Lipschitz inverse of $\Phi$.  Hence
\[
 \Hh^d\bigl(\{t\in C_d:H_t\cap V_s\cap Z\ne\varnothing\}\bigr)=0.
\]
The same argument with the coordinates exchanged handles a locally null
$H_t\cap Z$.  Thus all hypotheses of De Pauw's Theorem~8.3 are satisfied,
and input \ref{input:vh} proves the claim.
\end{proof}

\section*{Restriction to a measurable subspace}

\begin{lemma}[Trace and principal ideal]\label{lem:trace}
Let $\mu=\Hh^d$ on $[0,1]$, let $K\subset[0,1]$ be Borel, and let
$\phi=\mu|_K$ as an outer measure on the metric subspace $K$.  Then
\begin{align*}
 \A_\phi&=\{A\cap K:A\in\A_\mu\},\\
 \N_\phi^{\mathrm{loc}}&=
 \{N\cap K:N\in\N_\mu^{\mathrm{loc}}\}.
\end{align*}
Consequently
$\A_\phi/\N_\phi^{\mathrm{loc}}$ is isomorphic to the principal ideal below
$[K]$ in $\A_\mu/\N_\mu^{\mathrm{loc}}$.
\end{lemma}

\begin{proof}
Hausdorff outer measure on the metric subspace agrees with ambient Hausdorff
outer measure on subsets of $K$.  The trace of an ambient measurable set is
therefore $\phi$-measurable.  Conversely, suppose $E\subset K$ is
$\phi$-measurable.  Since the Borel set $K$ is $\mu$-measurable, for arbitrary
$S\subset[0,1]$ we have
\begin{align*}
 \mu(S)
 &=\mu(S\cap K)+\mu(S\setminus K)\\
 &=\mu(S\cap E)+\mu((S\cap K)\setminus E)+\mu(S\setminus K)\\
 &=\mu(S\cap E)+\mu(S\setminus E).
\end{align*}
Thus $E\in\A_\mu$, proving the first identity.

If $N\subset K$ is locally null in the ambient space, it is locally null in
$K$.  Conversely, if it is locally null in $K$ and $F\in\A_\mu$ has finite
measure, then $F\cap K$ is a finite-measure member of $\A_\phi$ and
\[
 \mu(N\cap F)=\phi(N\cap(F\cap K))=0.
\]
This proves the second identity.  The map $[E]\mapsto[E]$ identifies the
subspace quotient with the classes below $[K]$: if $[A]\leq[K]$, then
$A\setminus K$ is locally null and $[A]=[A\cap K]$.  Hence the image is
exactly the principal ideal below $[K]$.
\end{proof}

\section*{Lower-half resolution}

\begin{theorem}\label{thm:main}
For every $0<d<\frac12$, the statement
\[
 ([0,1],\A_{\Hh^d},\Hh^d)\text{ is semilocalizable}
\]
is undecidable in ZFC, in the usual relative-consistency sense.  More
precisely:
\begin{enumerate}[label=\textup{(\roman*)}]
\item under the Continuum Hypothesis the space is semilocalizable;
\item under
$\operatorname{non}(\N_{\mathcal L^1})<
 \operatorname{cov}(\N_{\mathcal L^1})$
it is not semilocalizable.
\end{enumerate}
\end{theorem}

\begin{proof}
Part (i) is input \ref{input:ch}, applied to the Borel-regular outer measure
$\Hh^d$ on the Polish space $[0,1]$.

For (ii), let $K_d$ be the compact set constructed above.  Lemma~\ref{lem:vh}
says that the quotient associated with $K_d$ is not order complete.  If the
interval were semilocalizable, then
\[
 \A_{\Hh^d}([0,1]) / \N_{\Hh^d}^{\mathrm{loc}}([0,1])
\]
would be order complete.  Every principal ideal in an order-complete Boolean
algebra is order complete: the supremum of a family below $b$ is the ambient
supremum, which is still below $b$.  Lemma~\ref{lem:trace} identifies the
quotient for $K_d$ with the principal ideal below $[K_d]$, giving a
contradiction.  Therefore the interval is not semilocalizable under the
stated consistent inequality.
\end{proof}

\begin{corollary}
Question~6 has an affirmative answer for every $0<d<\frac12$.  Together with
De Pauw's Theorem~9.8 at $d=\frac12$, undecidable semilocalizability is known
throughout $0<d\leq\frac12$.
\end{corollary}

\section*{Why the same route stops above one half}

The construction has a sharp geometric barrier.  Four separated
first-generation intervals of ratio $r=2^{-1/d}$ fit in an interval only if
$4r\leq1$, equivalently $d\leq\frac12$.  More invariantly,
\[
 \dim_H(C_d\times C_d)=2d>1\qquad(d>\tfrac12),
\]
so no bi-Lipschitz copy of the product can lie in $\R$.  This does not show
that Q6 is false in the upper half.  It shows only that any upper-half proof
must replace the equal-dimensional product-incidence geometry by a different
mechanism.

\section*{Novelty check and review status}

A bounded search of the run indexes, exact-title web/arXiv results, the source
paper's publication record, and later work on localizable locally determined
measurable spaces found no later paper explicitly answering Q6.  The proof is
short because its new ingredient is an elementary geometric transport of the
source paper's own Theorem~8.3.  Novelty remains provisional until specialist
literature review.

The main verifier points are the uniform first-different-digit estimates,
the passage from local nullity to nullity on finite-measure leaves, and the
trace/principal-ideal identification.  Each is proved explicitly above.

\begin{thebibliography}{9}
\bibitem{DePauw}
T.~De Pauw,
\emph{Undecidably semilocalizable metric measure spaces},
Commun. Contemp. Math. \textbf{26} (2024), no.~4, 2350012;
arXiv:1909.10190,
\url{https://arxiv.org/abs/1909.10190}.
\end{thebibliography}

\end{document}
